\clearpage\SourceLeaf{22}
\begin{center}
\bfseries ZUR THEORIE DER KUBISCHEN RESTE.
\end{center}



\begin{center}
[1.]
\end{center}

\[ p=mm+3nn \text{ eine Primzahl.} \]

\noindent\hspace*{\parindent}Für eine beliebige, durch $p$ nicht teilbare ganze Zahl $x$ setze man
\[
\left[\dfrac{2mx}{p}\right]=\alpha,\qquad \left[\dfrac{(3n-m)x}{p}\right]=\beta,\qquad \left[\dfrac{(-3n-m)x}{p}\right]=\gamma.
\]
\noindent Hieraus ergibt sich, wie man leicht sieht,
\[
\dfrac{2mx}{p}-\alpha+\dfrac{(3n-m)x}{p}-\beta+\dfrac{(-3n-m)x}{p}-\gamma
\]
als Summe dreier echter Brüche. Sie muss daher strikt zwischen den
Grenzen 0 und 3 liegen. Da diese Summe aber offenbar $=-\alpha-\beta-\gamma$, also
eine ganze Zahl ist, kann sie nur die Werte $+1$ und $+2$ annehmen.
Es gilt also entweder
\[ \alpha+\beta+\gamma=-1\quad\text{oder}\quad \alpha+\beta+\gamma=-2. \]

\noindent Hieraus folgt ferner, dass die Größen $\alpha, \beta, \gamma$ nicht alle drei zueinander
modulo 3 kongruent sein können. Unter den Differenzen $\alpha-\beta$, $\beta-\gamma$,
$\gamma-\alpha$ wird also mindestens eine (und eben deshalb mindestens noch eine,
